\chapter*{Abstract}
\addcontentsline{toc}{chapter}{Abstract}

This dissertation describes the design, implementation, and evaluation of an
\textit{Investor Sentiment Dashboard}: an end-to-end analytics system that applies
Natural Language Processing (NLP) and Explainable Artificial Intelligence (XAI)
to the multi-source analysis of public investor sentiment and its relationship to
financial asset price movements.

Three novel contributions are presented, each supported by empirical or
demonstrable results.
First, a \textit{finance-specific emotion taxonomy} layers six domain emotions
(fear, optimism, uncertainty, confidence, scepticism, and mixed) over FinBERT's
softmax probabilities using normalised Shannon entropy, domain lexicon matching,
and financial aspect priors.
No prior system in the reviewed literature combines these signals within a single,
architecture-agnostic scoring function; the module is deployed in production and
annotates the full 150{,}000-record corpus.
Second, a lightweight \textit{Aspect-Based Sentiment Analysis (ABSA)} module
extracts and tags evidence sentences across six financial aspect categories without
requiring fine-tuned sequence models, providing traceable justification for each
headline sentiment label.
Third, \textit{production-grade XAI integration}: LIME \citep{ribeiro2016lime} is
wired through the full application stack as a REST endpoint and rendered as a
token-level heatmap within the React SPA, making this one of the first reported
financial sentiment systems to surface token-level explanations directly in a
deployed user interface rather than an offline notebook.
The LIME integration also yielded a concrete diagnostic finding: contextually neutral
risk vocabulary (e.g.\ ``cautious'', ``inflation'') is systematically pushed toward
the negative class, directly identifying the mechanism behind FinBERT's neutral
recall deficit and pointing to a targeted fine-tuning strategy.

To underpin these contributions, \textbf{FinBERT} \citep{araci2019finbert} is
deployed as the core classifier on text aggregated from three distinct channels:
Reddit (via PRAW), X/Twitter (via a resilient multi-backend layer), and financial
news (via NewsAPI).
On a purpose-built, 200-sample evaluation benchmark reflecting the real
distribution of multi-source financial text, FinBERT achieves \textbf{82.0\%}
accuracy and \textbf{80.3\%} macro F1-score, \textbf{outperforming a
keyword-based baseline by 13.5 and 12.4 percentage points respectively}.
The most striking individual gain is on the negative class, where FinBERT recall
improves from 47.1\% to 94.3\%, a 47.2 percentage-point improvement that
illustrates the inherent limitation of lexicon-based methods on implicit financial
phrasing.

In addition to the three primary contributions, a \textit{sentiment--price
correlation engine} provides Pearson and Spearman correlation, configurable lag
analysis, Granger causality testing \citep{granger1969}, rolling correlation, and
market-adjusted metrics within a single interactive interface; a combination of
statistical methods not previously described together in an integrated,
ticker-agnostic dashboard.

Overall, the project demonstrates that multi-source, domain-adaptive sentiment
analysis can be deployed transparently and at scale within a single academic
system, and that the finance-emotion taxonomy, ABSA enrichment layer, and
production XAI integration in particular offer a publishable contribution at the
intersection of NLP, behavioural finance, and explainable AI.
A live industry demonstration of the delivered system to senior finance and AI
practitioners has been arranged for the summer 2026 period, with a view to
obtaining structured practitioner feedback and exploring potential real-world
adoption pathways beyond the scope of this dissertation.
